\section{空间群与 Bloch 定理}
\label{chap:space-groups-bloch}

点群只问哪些操作把一个有限物体留在原处。晶体还允许把整个结构平移一个晶格矢量，因此对称群变成无限离散群，并可能出现螺旋轴和滑移面这类点操作与分数平移的组合。

平移子群是正规子群，它的一维表示直接给出 Bloch 相位。固定一个波矢后，并非所有空间群元素都仍把它留在同一倒格矢等价类；剩下的元素组成小群，并决定该波矢处能带允许的简并和表示标签。

\subsection{晶体点群与空间群}
\label{sec:crystallographic-point-space}

普通分子只要求在有限范围内保持几何形状，晶体则多出一个全局条件：结构在三个独立方向上周期重复。这个平移周期性会反过来限制允许的旋转。最著名的结果就是晶体制约定理：周期晶格不允许五重旋转，也不允许一般的七重、八重旋转。下面直接由晶格基底下的整数矩阵与旋转迹推出只允许 $1,2,3,4,6$ 重旋转轴。

\begin{definitionplain}[Bravais 晶格]
取三个线性无关的基矢 $\vb a_1,\vb a_2,\vb a_3$。集合
\[
\mathcal L=
\left\{
 n_1\vb a_1+n_2\vb a_2+n_3\vb a_3
 \mid n_i\in\mathbb Z
\right\}
\]
称为三维 Bravais 晶格。若点群操作 $R$ 满足 $R\mathcal L=\mathcal L$，则称 $R$ 为晶格的点对称操作。
\label{def:bravais-lattice}
\end{definitionplain}

如果 $R$ 保持晶格，则每个基矢的像仍是晶格矢量，因此存在整数矩阵 $M(R)\in\operatorname{GL}(3,\mathbb Z)$，使
\[
R\vb a_j=\sum_{i=1}^3 M_{ij}(R)\vb a_i.
\label{eq:lattice-integer-representation}
\]
这个矩阵其实就是同一个线性变换在晶格基底下的表示。它与正交基下的旋转矩阵相似，因此具有相同的迹。

\begin{theorem}[晶体制约定理]
若周期晶格允许有限阶纯转动 $R$，设其转角为 $\theta=2\pi/n$，则
\[
n\in\{1,2,3,4,6\}.
\]
因此三维周期晶体不允许五重或一般 $n\ge7$ 的有限旋转轴。
\label{thm:crystallographic-restriction}
\end{theorem}

\par\medskip
\noindent\textbf{晶体制约定理}\quad
保持晶格的转动在晶格基底下由 $M\in\operatorname{GL}(3,\ZZ)$ 表示，所以 $\tr M\in\ZZ$。它与正交基底下的 $R$ 相似，而三维转动的本征值为 $1,\ee^{\ii\theta},\ee^{-\ii\theta}$，故
\[
\tr M=\tr R=1+2\cos\theta\in\ZZ.
\]
又因 $-1\leq\cos\theta\leq1$，只能有
\[
\cos\theta\in\left\{-1,-\frac12,0,\frac12,1\right\}.
\]
对有限阶转动取 $0\leq\theta\leq\pi$，相应转角为 $\pi,2\pi/3,\pi/2,\pi/3,0$，也就是 $n=2,3,4,6,1$。
\par\medskip

证明的关键是周期性带来的整数迹，而不是正五边形不能铺满平面的图形直觉。准晶可以具有五重或十重对称，正是因为它没有定理所要求的三维周期平移群；分子也没有这一限制，所以 $I_h$ 可用于描述富勒烯等有限结构。

晶体制约定理排除了五重轴等非晶体学旋转。把允许的旋转阶数 $1,2,3,4,6$ 系统地组合，就得到 32 个晶体点群和 7 个晶系。

“分子点群”允许任意有限阶转轴，但晶体不同。晶体包含平移周期性，一个点群操作不仅要保持局域几何，还必须把整套晶格点重新映回晶格点。这个额外要求会排除五重轴等许多几何上完全合法的旋转。所谓“32 个晶体点群”，就是所有与三维晶格兼容的有限点群，而不是人为选出的 32 个常用例子。

七个晶系则是进一步按最高阶旋转轴和整体轴系结构来组织这 32 个群。初学时不要把“晶系”“点群”“Bravais 晶格”当成同义词：点群只描述保持某一点不动的转动/反射部分；Bravais 晶格描述纯平移周期；晶系则是根据点群对晶格度量的约束形成的分类。把 \Cref{thm:crystallographic-restriction} 应用于上一节的 Schoenflies 家族，就得到三维晶体允许的 $32$ 个点群。它们进一步按照宏观轴系的共同特征归入 $7$ 个晶系。与其把 $32$ 个名字看成一张要背的表，不如把晶系理解成“晶胞度量约束 + 最高阶轴类型”的联合分类。

\begin{table}[htbp]
\centering
\caption{七个晶系、常用晶胞约束与 32 个晶体点群}
\label{tab:32-crystal-point-groups}
\small
\begin{tabularx}{\textwidth}{>{\bfseries}l l X}
\toprule
晶系 & 常用晶胞约束 & Schoenflies 点群 \\
\midrule
三斜 & $a\neq b\neq c$，角度一般 & $C_1,C_i$ \\
单斜 & $a\neq b\neq c$，两角为 $90^\circ$ & $C_2,C_s,C_{2h}$ \\
正交 & $a\neq b\neq c$，$\alpha=\beta=\gamma=90^\circ$ & $D_2,C_{2v},D_{2h}$ \\
四方 & $a=b\neq c$，三角皆 $90^\circ$ & $C_4,S_4,C_{4h},D_4,C_{4v},D_{2d},D_{4h}$ \\
三方 & 菱方或六角设置 & $C_3,S_6,D_3,C_{3v},D_{3d}$ \\
六方 & $a=b\neq c$，$\gamma=120^\circ$ & $C_6,C_{3h},C_{6h},D_6,C_{6v},D_{3h},D_{6h}$ \\
立方 & $a=b=c$，三角皆 $90^\circ$ & $T,T_h,O,T_d,O_h$ \\
\bottomrule
\end{tabularx}
\end{table}

例如，四方晶系的核心不是“名字里出现 4”，而是存在唯一的四重主轴，同时基面两个方向等价而与主轴方向不等价；这正对应 $a=b\neq c$。立方晶系则有多条三重轴与四重/二重轴，使三个笛卡尔方向在宏观上完全等价，因此 $a=b=c$。这种几何理解在判断张量物性的独立分量时比死记点群名字更有用。

只知道允许的旋转阶数还不够；晶体点群还必须保持具体晶格的度量结构。把这一点写成矩阵方程后，七个晶系的来源会更清楚。

晶胞参数 $a,b,c,\alpha,\beta,\gamma$ 可以统一编码进晶格基底的 Gram 矩阵
\[
G_{ij}=\vb a_i\cdot\vb a_j.
\]
如果晶格点操作在这组基下由整数矩阵 $M$ 表示，那么保持内积等价于
\[
M^{\mathsf T}GM=G.
\label{eq:lattice-metric-invariance}
\]
因此一个给定晶格的点群可以理解为满足 \eqnrefs{eq:lattice-metric-invariance} 的有限整数矩阵群。这个表达把“晶胞长什么样”与“允许哪些点群操作”直接联系起来。

例如立方晶格有
\[
G=a^2 I_3,
\]
因此三个坐标方向完全等价，允许大量置换和符号翻转；四方晶格则有
\[
G=\operatorname{diag}(a^2,a^2,c^2),
\qquad a\neq c,
\]
只有 $x,y$ 两方向等价，$z$ 方向被区分出来，因此最多保留唯一四重主轴。正交晶格的
\[
G=\operatorname{diag}(a^2,b^2,c^2),
\qquad a,b,c\ \text{互异},
\]
则进一步把三方向全部区分，只留下相互垂直的二重轴结构。这样，七晶系的参数限制就不再只是表格，而是晶格度量张量稳定子的不同可能。

这里尤其要解释一个容易混淆的问题：三方与六方晶系都常用六方坐标描述，但它们不是同一个晶系。

三方和六方经常让初学者困惑，因为它们都可以采用六角轴系来画晶胞。关键区别不是“晶胞画法”，而是最高允许旋转阶数。六方晶系存在六重轴；三方晶系的最高主轴是三重轴。三方晶体可以用菱方原胞，也可以嵌入六角常规晶胞描述，因此在晶胞参数层面会出现两套常见设置。真正不变的是点群与晶格的对称结构，而不是某一种画法。

点群分类完成以后，还要把平移骨架的可能性组织起来，这就引出十四种 Bravais 晶格。


点群描述“允许哪些旋转”，晶格描述“平移骨架怎样重复”。同一个晶系还可能有不同的中心化方式。将七个晶系与允许的中心化方式结合，得到十四种 Bravais 晶格。

\begin{table}[htbp]
\centering
\caption{三维十四种 Bravais 晶格}
\label{tab:14-bravais}
\begin{tabular}{c|c}
\toprule
晶系 & Bravais 晶格类型 \\
\midrule
三斜 & $P$ \\
单斜 & $P,C$ \\
正交 & $P,C,I,F$ \\
四方 & $P,I$ \\
三方 & $R$ \\
六方 & $P$ \\
立方 & $P,I,F$ \\
\bottomrule
\end{tabular}
\end{table}

这里 $P$、$I$、$F$、$C$、$R$ 分别表示原始、体心、面心、底心和菱方中心化。点群与 Bravais 晶格并不是任意配对；它们必须共享同一晶系的度量结构。空间群正是进一步把“点操作”和“平移”组合起来的完整群。

十四种 Bravais 晶格来自两个条件：中心化点集必须在晶体点群下保持不变，而由幺模整数基变换联系的晶胞描述同一晶格。因此某些看似新的中心化会在换基后退化为表中的类型；例如单斜的体心描述可换成底心描述，四方的面心描述可换成体心描述。这个等价关系解释了为何“七晶系乘若干中心化”的朴素计数会重复。

惯用晶胞优先显露对称性，不一定是最小原胞。能带和态密度计算则必须用原胞体积 $\Omega$；相应的倒格原胞体积为 $(2\pi)^3/\Omega$。在实空间取 Wigner--Seitz 原胞，在倒空间便得到第一 Brillouin 区。

有了平移晶格的分类，下一步自然转向其对偶结构——倒格子。它把晶面、衍射条件与 Bloch 波矢统一到同一套语言中。

晶面在三个晶格方向上的截距倒数化成互素整数 $(h,k,l)$，便得到 Miller 指数；平行于某个基矢的晶面在该方向的指数为零。定义倒格基矢满足 $\vb a_i\cdot\vb b_j=2\pi\delta_{ij}$ 后，每个晶面族对应
\[
\vb G_{hkl}=h\vb b_1+k\vb b_2+l\vb b_3,
\]
其中 $\vb G_{hkl}$ 垂直于晶面，且相邻晶面间距为
\[
d_{hkl}=\frac{2\pi}{|\vb G_{hkl}|}.
\]

\par\medskip
\noindent\textbf{晶面法向与面间距}\quad
第 $n$ 张平行晶面上的点 $\vb r=x_1\vb a_1+x_2\vb a_2+x_3\vb a_3$ 满足
\[
hx_1+kx_2+lx_3=n,
\qquad n\in\ZZ.
\]
任取同一晶面内两点之差 $\Delta\vb r$，上式给出 $\vb G_{hkl}\cdot\Delta\vb r=0$，故 $\vb G_{hkl}$ 是法向。相邻两面的位移 $\Delta\vb r$ 满足 $\vb G_{hkl}\cdot\Delta\vb r=2\pi$；沿单位法向取最短位移便有
\[
d_{hkl}|\vb G_{hkl}|=2\pi.
\]
\par\medskip

弹性散射保持晶格平移的特征标，因而要求 $\ee^{\ii(\vb k'-\vb k)\cdot\vb R}=1$ 对所有 $\vb R\in\mathcal L$ 成立，即 Laue 条件 $\vb k'-\vb k\in\mathcal L^*$。衍射峰落在倒格点上，不是额外的经验规则，而是平移群特征标的正交性。

晶格的 Fourier 变换是倒格子上的 Dirac 梳：
\[
\hat\Lambda(\vb k)
=\sum_{\vb R\in\Lambda}\ee^{-i\vb k\cdot\vb R}
=\frac{(2\pi)^3}{\Omega}
\sum_{\vb G\in\Lambda^*}\delta(\vb k-\vb G).
\]
这一 Poisson 求和公式把实空间周期性转化为倒空间离散性，是衍射物理的数学基础。结构因子是电子密度的 Fourier 系数，
\[
F(\vb G)=\int\rho(\vb r)\ee^{-i\vb G\cdot\vb r}\,\dd^3r,
\]
因此衍射实验测量的是 $\rho$ 的 Fourier 变换在倒格点上的值。

设 $A=[\vb a_1|\vb a_2|\vb a_3]$ 与 $B=[\vb b_1|\vb b_2|\vb b_3]$。由 $A^{\mathsf T}B=2\pi I_3$ 取行列式，立刻得到实、倒格原胞体积的关系
\[
\Omega\,\Omega^*=(2\pi)^3.
\]
Born--von Karman 边界条件下，$N$ 个原胞给出 $N$ 个允许波矢，每个波矢在 Brillouin 区占体积 $\Omega^*/N$，所以
\[
\frac{1}{N}\sum_{\vb k}\to\frac{\Omega}{(2\pi)^3}\int_{\mathrm{BZ}}\dd^3k,
\]
这也是能带积分中测度因子的来源。Bloch 波矢只需取第一 Brillouin 区，因为 $\vb k$ 与 $\vb k+\vb G$ 给出同一平移特征标。

倒格子揭示了晶体的平移对称性，而点群对称性同样对宏观物理性质施加约束。晶体点群不仅约束能级结构，还直接限制物理性质张量的独立分量个数——这是群论在材料科学中最实用的结论之一。

\begin{theorem}[张量不变性原理]
设晶体点群为 $P$。物理性质张量 $T_{i_1\cdots i_n}$（$n$ 阶）在 $P$ 的所有操作 $R$ 下不变：
\[
T_{i_1\cdots i_n}
=\sum_{j_1,\dots,j_n}
R_{i_1 j_1}\cdots R_{i_n j_n}\,T_{j_1\cdots j_n},
\qquad\forall R\in P.
\]
此方程组确定 $T$ 的独立分量个数。
\label{thm:tensor-invariance}
\end{theorem}

二阶对称张量（如介电张量 $\epsilon_{ij}$、电导率 $\sigma_{ij}$）在转动 $R$ 下变换为 $\epsilon'=R\epsilon R^{\mathsf T}$。不变性要求 $R\epsilon=\epsilon R$，即 $\epsilon$ 与所有点群操作对易。等价地，$\epsilon$ 属于算符表示空间中的全对称表示 $A_1$ 分量。独立分量数等于
\[
N_{A_1}=\frac{1}{|G|}\sum_{g\in G}\chi_{\mathrm{Sym}^2}(g),
\]
其中 $\chi_{\mathrm{Sym}^2}(g)$ 是二阶对称张量表示的特征标。对三维矢量表示 $R(g)$，$\chi_{\mathrm{Sym}^2}(g)=\frac12[(\tr R)^2+\tr(R^2)]$。

对称性越高，不变张量空间通常越小。二阶对称张量在三斜、单斜、正交、四方/三方/六方和立方晶系中分别有 $6,4,3,2,1$ 个独立分量。因而单轴晶体有 $\epsilon=\operatorname{diag}(\epsilon_\perp,\epsilon_\perp,\epsilon_\parallel)$，立方晶体则只能有 $\epsilon=\epsilon_0 I_3$。四阶弹性张量还受指标交换对称约束；立方晶体最终只留下 $C_{11},C_{12},C_{44}$ 三个独立常数。

这些特例可统一为张量表示中的不动子空间问题。

设 $V=\RR^3$ 是点群 $P$ 的自然表示。$n$ 阶张量属于 $V^{\otimes n}$，群元 $R$ 的作用为
\[
(R\cdot T)_{i_1\cdots i_n}=\sum_{j_1,\dots,j_n}R_{i_1j_1}\cdots R_{i_nj_n}\,T_{j_1\cdots j_n}
\]
若张量具有指标交换对称性，应先把表示空间限制到相应的对称化子空间，再施加点群不变性。

\begin{theorem}[张量独立分量的计数]
\label{thm:tensor-count-general}
设 $P$ 是晶体点群，$T$ 是秩 $n$ 张量，其对称性由表示 $\rho_T:P\to\operatorname{GL}(W)$ 描述（$W\subseteq V^{\otimes n}$ 是对称性约束后的子空间）。则 $T$ 在 $P$ 对称性下的独立分量数为
\begin{equation}
N=\frac1{|P|}\sum_{R\in P}\chi_T(R),
\label{eq:tensor-count-general}
\end{equation}
其中 $\chi_T$ 是 $\rho_T$ 的特征标。等价地，$N$ 是平凡表示 $A_1$ 在 $\rho_T$ 中的重数。
\end{theorem}

群平均算符
\[
\Pi_P=\frac1{|P|}\sum_{R\in P}\rho_T(R)
\]
满足 $\Pi_P^2=\Pi_P$，像空间正是不变张量空间；取迹便得到 \eqnrefs{eq:tensor-count-general}。这段论证也说明，公式计算的是已有指标对称约束之后的独立分量，而不是未经约束的 $3^n$ 个分量。

\begin{proposition}[对称幂表示的特征标]
\label{prop:sym-power-char}
设 $R\in P$ 有本征值 $\lambda_1,\lambda_2,\lambda_3$（作为 $3\times3$ 正交矩阵）。则：
\begin{enumerate}[label=(\roman*),leftmargin=3em]
\item 全张量空间 $V^{\otimes n}$ 的特征标：$\chi_{V^{\otimes n}}(R)=(\lambda_1+\lambda_2+\lambda_3)^n=(\tr R)^n$。
\item 对称幂 $\mathrm{Sym}^n V$ 的特征标：
\begin{equation}
\chi_{\mathrm{Sym}^n}(R)=\sum_{\substack{k_1+k_2+k_3=n\\k_i\ge0}}\lambda_1^{k_1}\lambda_2^{k_2}\lambda_3^{k_3},
\label{eq:sym-power-char}
\end{equation}
即 $(\lambda_1,\lambda_2,\lambda_3)$ 的次数 $n$ 完全齐次对称多项式。
\item 反对称幂 $\wedge^n V$ 的特征标：$\chi_{\wedge^n}(R)=\sum_{i_1<\cdots<i_n}\lambda_{i_1}\cdots\lambda_{i_n}$（初等对称多项式 $e_n$）。
\end{enumerate}
\end{proposition}

低阶计算最常用的特例是
\[
\chi_{\mathrm{Sym}^2}(R)
=\frac12\bigl[\chi_V(R)^2+\chi_V(R^2)\bigr].
\]
\begin{proposition}[张量表示的不可约分解]
\label{prop:tensor-irrep-decomposition}
张量表示分解为
\[
\rho_T\simeq\bigoplus_\Gamma\,m_\Gamma\,\Gamma,\qquad m_\Gamma=\frac1{|P|}\sum_{R\in P}\overline{\chi_\Gamma(R)}\,\chi_T(R).
\]
平凡表示的重数 $m_{A_1}$ 给出不变张量数，其他分量则描述在外场或相变中按相应不可约表示变化的响应模式。
\end{proposition}

32 个点群分类了晶体的旋转对称性，而空间群进一步把点操作与平移组合，产生 Wyckoff 位置——这是描述晶体内原子占位对称性的基本工具。

\begin{definitionplain}[Wyckoff 位置]
空间群 $\mathcal G$ 作用下，空间中一般位置 $\vb r$ 的轨道 $|\mathcal G\cdot\vb r|$ 大小等于 $|\mathcal G/T|$（点群阶数）。若 $\vb r$ 位于对称元素上，其稳定子 $G_{\vb r}$ 非平凡，轨道变小。具有相同稳定子共轭类的位置集合称为一个 Wyckoff 位置，用字母 $a,b,c,\dots$ 标记（多重度递增）。
\end{definitionplain}

\medskip
\noindent\emph{空间群 $Pm\bar3m$（$O_h^1$，No.221）的 Wyckoff 位置。}
立方 $Pm\bar3m$ 的部分 Wyckoff 位置：
\[
\begin{array}{c|c|c|l}
\text{Wyckoff} & \text{多重度} & \text{点对称} & \text{坐标}\\\hline
a&1&m\bar3m&(0,0,0)\\
b&1&m\bar3m&(\tfrac12,\tfrac12,\tfrac12)\\
c&3&4mm&(0,\tfrac12,\tfrac12)\\
d&6&mmm&(\tfrac12,0,0)\\
e&12&\bar42m&(\tfrac14,\tfrac14,0)
\end{array}
\]
位置 $a$（角点）和 $b$（体心）具有完整 $O_h$ 对称，适合放完全对称的原子（如 CsCl 中 Cs 在 $a$、Cl 在 $b$）。位置 $c$（面心）有 $C_{4v}$ 对称，适合放面心原子。Wyckoff 位置直接决定原子在晶体中的对称环境和可能的局部模式。
\par\medskip

Wyckoff 位置描述了原子在空间群作用下的对称环境，而空间群本身的结构则需要用 Seitz 记号来系统刻画。

点群只记录 $\vb r\mapsto R\vb r$。晶体的完整对称操作还允许在转动之后附加平移，因此一般形式必须写成
\[
\vb r\longmapsto R\vb r+\vb t.
\]
Seitz 记号 $\{R\mid\vb t\}$ 只是把这一仿射变换紧凑地记下来。竖线左边是线性部分 $R$，右边是平移部分 $\vb t$。特别地，$\{I\mid\vb a\}$ 是纯平移，而 $\{R\mid\vb 0\}$ 是过原点的纯点群操作。

乘法公式不能凭记忆写。若先做 $\{S\mid\vb s\}$ 再做 $\{R\mid\vb t\}$，位置依次变为
\[
\vb r\xmapsto{\{S\mid\vb s\}}S\vb r+\vb s
\xmapsto{\{R\mid\vb t\}}R(S\vb r+\vb s)+\vb t
=RS\vb r+(R\vb s+\vb t),
\]
所以
\[
\{R\mid\vb t\}\{S\mid\vb s\}=\{RS\mid R\vb s+\vb t\}.
\]
注意第二个平移 $\vb s$ 会先被前面的转动 $R$ 旋转，这正是空间群通常不是点群与平移群简单直积的原因。一个空间群操作写成 Seitz 记号
\[
\{R\mid\vb t\}:\vb r\longmapsto R\vb r+\vb t.
\]
这里 $R\in O(3)$ 是点操作，$\vb t$ 是平移。重要的是，$\vb t$ 不一定本身就是晶格矢量；在螺旋轴和滑移面中，它可以是晶格周期的分数倍。

\begin{proposition}[Seitz 乘法与逆元]
空间操作满足
\[
\{R\mid\vb t\}\{S\mid\vb s\}
=\{RS\mid R\vb s+\vb t\},
\label{eq:seitz-product}
\]
以及
\[
\{R\mid\vb t\}^{-1}
=\{R^{-1}\mid- R^{-1}\vb t\}.
\label{eq:seitz-inverse}
\]
\label{prop:seitz-algebra}

\end{proposition}

\begin{proof}
连续作用到任意点 $\vb r$：
\begin{align*}
\{R\mid\vb t\}\{S\mid\vb s\}\vb r
&=\{R\mid\vb t\}(S\vb r+\vb s)\\
&=RS\vb r+R\vb s+\vb t,
\end{align*}
因此得到 \eqnrefs{eq:seitz-product}。要求逆操作满足
\[
\{R\mid\vb t\}\{R^{-1}\mid\vb u\}=
\{I_3\mid\vb0\},
\]
由乘法式知 $R\vb u+\vb t=0$，故 $\vb u=-R^{-1}\vb t$，得到 \eqnrefs{eq:seitz-inverse}。
\end{proof}

所有纯平移
\[
T=\{\{I_3\mid\vb R\}\mid \vb R\in\mathcal L\}
\]
构成空间群 $\mathcal G$ 的正规 Abel 子群。把平移部分“商掉”以后，剩下的就是有限的晶体点群：
\[
\mathcal G/T\simeq P,
\]
其中 $P$ 是空间群的点群。这一短正合结构
\[
1\longrightarrow T\longrightarrow\mathcal G
\longrightarrow P\longrightarrow1
\]
是理解空间群最干净的代数框架。\chapref{chap:group-theory-quantum}讲 Bloch 定理和小群时会不断使用它。

如果通过选择适当原点，可以让每个点群元素 $R$ 都在空间群中以 $\{R\mid0\}$ 出现，那么空间群可以看成平移群与点群的半直积，称为\textbf{对称型空间群}（symmorphic 空间群；传统教材也常称“简单空间群”）。三维共有 $73$ 个。若某些点操作只能与分数平移绑在一起存在，就得到\textbf{非对称型空间群}（nonsymmorphic 空间群；传统教材也常称“非简单空间群”），共有 $157$ 个。两类合计 $230$ 个空间群。

非对称型空间群最典型的几何元素是螺旋轴与滑移面。它们的关键不是“多做了一次平移”，而是点操作 $R$ 与分数平移 $\vb t$ \emph{不可拆开分别成为对称操作}。

螺旋轴与滑移面的几何如图 \figref{fig:screw-glide} 所示。
\begin{figure}[htbp]
\centering
\begin{minipage}{0.47\textwidth}
\centering
\begin{tikzfig}[scale=0.86]
  \draw[axes] (0,-2.5)--(0,2.8) node[above] {$z$};
  \foreach \k in {0,...,5}{
    \pgfmathsetmacro{\ang}{60*\k}
    \pgfmathsetmacro{\yy}{-2+0.72*\k}
    \pgfmathsetmacro{\xx}{0.95*cos(\ang)}
    \fill (\xx,\yy) circle (1.5pt);
  }
  \draw[lecturearrow] (1.1,-1.9) to[out=25,in=-35] (0.65,-.9);
  \node[lecturelabel] at (2.05,-1.05) {$\{C_6\mid c/6\}$};
\end{tikzfig}
\small screw axis
\end{minipage}\hfill
\begin{minipage}{0.47\textwidth}
\centering
\begin{tikzfig}[scale=0.86]
  \draw[dashed,line width=.75pt] (-2.6,0)--(2.6,0) node[right,lecturelabel] {$\sigma$};
  \foreach \x in {-2,-.8,.4,1.6}{\fill (\x,1.0) circle (1.5pt);}
  \foreach \x in {-1.4,-.2,1.0,2.2}{\fill (\x,-1.0) circle (1.5pt);}
  \draw[lecturearrow] (-1.9,.75)--(-1.42,-.75);
  \node[lecturelabel] at (.25,.55) {$\{\sigma\mid \vb a/2\}$};
\end{tikzfig}
\small glide plane
\end{minipage}
\caption{非简单空间群的两种基本元素：螺旋轴与滑移面。点操作和分数平移必须作为一个整体出现。}
\label{fig:screw-glide}
\end{figure}

这类“不可拆分的分数平移”在纯几何分类中看似只是记号上的麻烦，但到了能带理论里会产生实质后果：在布里渊区边界，螺旋轴或滑移面的表示可能平方到一个非平凡的 Bloch 相位，从而强制某些能带粘连。\chapref{chap:group-theory-quantum}会把这个现象作为点群思维向空间群思维升级的例子。

Seitz 记号还揭示了一个重要事实：平移部分会随原点改变，因此判断一个空间群是否能“拆成点群与整格平移”必须考察哪些分数平移能够通过移动原点消去。

Seitz 记号中的平移部分依赖原点选择。若把原点从 $O$ 移到 $O'=O+\vb d$，同一个几何操作
\[
\{R\mid\boldsymbol\tau\}
\]
在新原点下的平移部分变为
\[
\boldsymbol\tau'
=\boldsymbol\tau+(R-I_3)\vb d.
\label{eq:origin-shift-seitz}
\]
因此看到一个分数平移时，不能立刻断言它是“真正的”螺旋或滑移成分；有些分数部分可以通过移动原点消掉。简单空间群之所以称为 symmorphic，就是因为存在某个原点使所有点操作都能选成 $\{R\mid0\}$ 的代表。

真正的非对称型空间群操作所含的本征分数平移则无法通过任何原点移动消去。例如绕轴的螺旋平移分量沿着旋转轴，而 $(R-I_3)\vb d$ 总垂直于固定轴方向，因此轴向分数平移不可能靠移原点去掉。同理，滑移面中平行镜面的分数平移也具有不可消除部分。这一几何事实正是“非简单空间群不是糟糕原点选择造成的”最直接的解释。

\medskip
\noindent\emph{一维链上的半格矢滑移类比。}
虽然真正的滑移面需要二维或三维空间，但一维类比很有助于理解。设两个子晶格 $A_n$ 与 $B_n$ 分别位于 $na$ 与 $na+a/2$。单独平移 $a/2$ 会交换 $A$、$B$，如果两者是不同原子就不是对称操作；但若同时伴随一个内部交换操作，组合可能恢复整个结构。空间群中的螺旋与滑移正是“点操作修正了分数平移造成的不匹配”的几何版本。
\par\medskip

进入晶体文献以后还会遇到另一套命名体系。Schoenflies 与 Hermann--Mauguin 符号并非互相竞争，而是承担不同的信息任务。

Schoenflies 符号特别适合分子光谱与点群表示，因为它突出主轴、镜面和多面体类型；晶体学更常用 Hermann--Mauguin 国际符号，因为后者同时编码晶格中心化以及沿特定晶向的旋转、螺旋和滑移信息。例如符号开头的 $P,I,F,C,R$ 直接告诉我们 Bravais 中心化类型，后续的数字与字母则描述相应方向上的轴或面。

因此，当文献给出一个空间群时，常见工作流是：先从 Hermann--Mauguin 符号或空间群编号确定完整的 $\{R\mid\boldsymbol\tau\}$ 操作；再把平移部分商掉得到点群；最后根据具体 $\vb k$ 取小群并查相应不可约表示。\chapref{chap:group-theory-quantum}的能带部分会沿这条链继续下去。

在实际识别点群时，三维对称元素常用极射赤面投影压缩到二维图上。


极射赤面投影并不引入新的群论定理，它的任务是把三维对称元素压缩到二维图上，使我们能够在同一张平面图上同时看见多个旋转轴与镜面。基本方法是：把旋转轴与单位球面的交点看成极点，再从球极向赤道平面作立体投影。上半球与下半球的点用不同符号区分。

立体投影把球面方向映射到赤道平面，如图 \figref{fig:stereographic-projection} 所示。
\begin{figure}[htbp]
\centering
\begin{tikzfig}[scale=1.0]
  \draw[line width=.75pt] (0,0) circle (2.2);
  \draw[dashed] (-2.2,0)--(2.2,0);
  \node[lecturelabel] at (-1.2,-0.62) {equator};
  \coordinate (N) at (0,2.2);
  \coordinate (S) at (0,-2.2);
  \coordinate (P) at (1.4,1.25);
  \fill (P) circle (1.5pt) node[above right,lecturelabel] {$P$};
  \coordinate (Pp) at (0.89,0);
  \draw[line width=.65pt] (S)--(P)--(Pp);
  \fill (Pp) circle (1.5pt);
  \node[lecturelabel] at (0.89,-0.62) {$P'$};
  \draw[line width=.6pt] (0,0) ellipse (2.2 and .38);
  \node[lecturelabel] at (-1.0,-2.65) {stereographic projection};
\end{tikzfig}
\caption{极射赤面投影的基本几何。实际晶体学图中会用约定符号标出不同阶数的轴、镜面与上下半球极点。}
\label{fig:stereographic-projection}
\end{figure}

对于本课程，掌握投影的几何思想和常见符号足够；真正分析量子态时，更重要的是从图中读出点群，再调用其不可约表示。于是本章最后一节转向这一层表示论结构。

\subsection{晶体点群的不可约表示}
\label{sec:crystal-point-group-irreps}

本节不把三十二张特征标表当成孤立数字表来背，而是解释这些表怎样由点群结构组织出来，以及为什么表中会出现 $A,B,E,T$、$g/u$ 这些标签。这样才能自然过渡到下一章的量子力学应用。

从点群的结构分类转入表示论：先从最简单的循环群和二面体群开始，它们是不可约表示的基石。

对循环群 $C_n=\langle r\mid r^n=e\rangle$，群是 Abel 的，因此所有不可约复表示都是一维。它们由生成元的像决定：
\[
D_m(r)=\exp\!\left(\frac{2\pi\ii m}{n}\right),
\qquad m=0,1,\dots,n-1.
\label{eq:cn-irreps}
\]
这已经给出全部 $n$ 个不可约表示。若希望使用实表示，则复共轭的一维表示往往会合并成二维实表示，这正是特征标表里 $E$ 型表示频繁出现的原因之一。

对 $D_n$，可取群表示
\[
D_n=\langle r,s\mid r^n=e,\ s^2=e,\ srs=r^{-1}\rangle.
\]
其中 $r$ 对应绕主轴的 $2\pi/n$ 转动，$s$ 对应一条横向二重轴。二维不可约表示可以写成
\[
D_m(r)=
\begin{pmatrix}
\cos\frac{2\pi m}{n}&-\sin\frac{2\pi m}{n}\\
\sin\frac{2\pi m}{n}&\cos\frac{2\pi m}{n}
\end{pmatrix},
\qquad
D_m(s)=
\begin{pmatrix}1&0\\0&-1\end{pmatrix}.
\]
因此
\[
\chi_m(r^k)=2\cos\frac{2\pi mk}{n},
\qquad
\chi_m(sr^k)=0.
\]
当 $n$ 为奇数时有两个一维不可约表示和 $(n-1)/2$ 个二维不可约表示；当 $n$ 为偶数时有四个一维不可约表示和 $n/2-1$ 个二维不可约表示。维数平方和分别验证
\[
2\cdot1^2+\frac{n-1}{2}\cdot2^2=2n,
\]
以及
\[
4\cdot1^2+\left(\frac n2-1\right)2^2=2n.
\]
这比逐个 $D_2,D_3,D_4,D_6$ 单独重算更有普适性。

循环群与二面体群之后，再看多面体群的不可约表示维数怎样由类数和平方和条件限制。


对 $T,O,I$，共轭类数和维数平方和已经把不可约表示的维数强烈限制住。标准结果为
\[
T:\quad 1,1,1,3,
\]
\[
O:\quad 1,1,2,3,3,
\]
\[
I:\quad 1,3,3,4,5.
\]
例如 $O$ 的群阶为 $24$，五个共轭类意味着五个不可约表示，而
\[
1^2+1^2+2^2+3^2+3^2=24.
\]
两个三维表示并不是“因为空间是三维所以必然有两个”，而是共轭类结构、正交关系与维数平方和共同决定的。

若群包含中心反演且分解成
\[
G=K\times C_i,
\]
则 $C_i$ 的两个一维不可约表示分别给出反演本征值 $+1$ 与 $-1$。于是 $K$ 的每个不可约表示都复制成一对：偶宇称记作 $g$（gerade），奇宇称记作 $u$（ungerade）。例如 $O_h=O\times C_i$，因此 $O$ 的五个不可约表示变成 $O_h$ 的十个不可约表示。

在分子与晶体文献中，最常见的是 Mulliken 符号。其结构应当按“维数 + 附加对称量子数”理解：$A$、$B$ 表示一维不可约表示，$E$ 表示二维，$T$（有些文献写 $F$）表示三维；下标 $1,2$ 区分同维表示在某些 $C_2$ 或镜面下的不同特征标；$g/u$ 标记反演偶奇；撇号 $' / ''$ 则常用来标记相对于水平镜面的偶奇。

因此，看到 $T_{2g}$ 时不要把它当作一个不可分解的名字。它至少告诉我们三件事：表示维数为 $3$；它属于第二种 $T$ 型；在中心反演下为偶。\chapref{chap:group-theory-quantum}讨论 $d$ 轨道的晶场劈裂时，这三个信息都会直接使用。

\begin{table}[htbp]
\centering
\caption{$C_{3v}$ 的特征标表与常见基函数}
\label{tab:c3v-character}
\begin{tabular}{c|ccc|l}
\toprule
 & $E$ & $2C_3$ & $3\sigma_v$ & 常见基函数 \\
\midrule
$A_1$ & $1$ & $1$ & $1$ & $z;\ x^2+y^2,z^2$ \\
$A_2$ & $1$ & $1$ & $-1$ & $R_z$ \\
$E$   & $2$ & $-1$ & $0$ & $(x,y);\ (R_x,R_y);\ (x^2-y^2,2xy)$ \\
\bottomrule
\end{tabular}
\end{table}

这张表与\chapref{chap:representation-theory} 中 $D_3$ 的特征标表\cref{tab:d3-character-table} 完全相同，只差记号：$E\leftrightarrow e$、$C_3\leftrightarrow r$、$\sigma_v\leftrightarrow s$。三个不可约表示 $A_1$、$A_2$、$E$ 分别对应 $D_3$ 的两个一维表示和二维表示。基函数列则给出了每个不可约表示在分子坐标上的具体载体——例如 $z$ 在 $A_1$ 中（$C_3$ 轴方向不变）、$(x,y)$ 在 $E$ 中（垂直于主轴的平面被 $C_3$ 转动），这些信息在\chapref{chap:group-theory-quantum} 构造对称性适配线性组合和判断光谱选择定则时将直接使用。

\begin{table}[htbp]
\centering
\caption{$C_{4v}$ 的特征标表与常见基函数}
\label{tab:c4v-character}
\begin{tabular}{c|ccccc|l}
\toprule
 & $E$ & $2C_4$ & $C_2$ & $2\sigma_v$ & $2\sigma_d$ & 常见基函数 \\
\midrule
$A_1$ & $1$&$1$&$1$&$1$&$1$ & $z,\ x^2+y^2,z^2$ \\
$A_2$ & $1$&$1$&$1$&$-1$&$-1$ & $R_z$ \\
$B_1$ & $1$&$-1$&$1$&$1$&$-1$ & $x^2-y^2$ \\
$B_2$ & $1$&$-1$&$1$&$-1$&$1$ & $xy$ \\
$E$   & $2$&$0$&$-2$&$0$&$0$ & $(x,y);\ (R_x,R_y);\ (xz,yz)$ \\
\bottomrule
\end{tabular}
\end{table}

这些“基函数”并不是表格的附属装饰，而是把表示论与物理变量连接起来的桥梁。例如 $(x,y,z)$ 告诉我们电偶极矩或位移矢量如何变换；$(x^2-y^2,xy,xz,yz,2z^2-x^2-y^2)$ 告诉我们 $d$ 轨道和 Raman 张量如何分解。下一章所有选择定则几乎都可以追溯到这些变换性质。

特征标表不应被当作只能查阅的黑箱。对许多几何表示，特征标可以直接从“不动对象贡献多少迹”算出来。

在具体点群中，很多表示并不需要逐个写矩阵才能求特征标。最常见的是三维矢量表示。若纯转动角为 $\theta$，由 \eqnrefs{eq:rotation-trace} 直接有
\[
\chi_{\mathrm{vec}}(R)=1+2\cos\theta.
\]
若是不保手性操作 $\mathcal I R$，三维矢量表示矩阵多一个整体负号，因此
\[
\chi_{\mathrm{vec}}(\mathcal I R)=-(1+2\cos\theta).
\]
这使得点群特征标表中的矢量基函数 $(x,y,z)$ 一栏可以从几何立刻验证。

对原子位移、等价轨道或等价原子构成的置换表示，还有另一条非常实用的规则：\emph{特征标等于在该操作下保持在原位置的基对象数，乘以每个对象内部自由度的迹。} 如果基对象只是标量轨道标签，置换表示的特征标就是固定点数；如果每个固定原子还带三维位移矢量，就再乘相应局域旋转矩阵的迹。\chapref{chap:group-theory-quantum}推导振动表示时会直接使用这一点。

有了置换表示的特征标计算规则，接下来以立方群 $O$ 为范例，看一张完整的特征标表怎样从群结构中自然生长出来，再通过直积扩展到 $O_h$。

立方群是凝聚态物理里最常见的高对称点群之一。$O$ 有五个共轭类，可以按
\[
E,\quad 8C_3,\quad 3C_2(C_4^2),\quad 6C_4,\quad 6C_2'
\]
排列。其五个不可约表示的维数为 $1,1,2,3,3$。一组常用特征标表为

\begin{table}[htbp]
\centering
\caption{纯立方旋转群 $O$ 的特征标表}
\label{tab:o-character}
\begin{tabular}{c|ccccc|l}
\toprule
 & $E$ & $8C_3$ & $3C_2$ & $6C_4$ & $6C_2'$ & 常见基函数 \\
\midrule
$A_1$ & $1$&$1$&$1$&$1$&$1$ & $x^2+y^2+z^2$ \\
$A_2$ & $1$&$1$&$1$&$-1$&$-1$ & $xyz(x^2-y^2)(y^2-z^2)(z^2-x^2)$ \\
$E$   & $2$&$-1$&$2$&$0$&$0$ & $(2z^2-x^2-y^2,\ x^2-y^2)$ \\
$T_1$ & $3$&$0$&$-1$&$1$&$-1$ & $(x,y,z)$ \\
$T_2$ & $3$&$0$&$-1$&$-1$&$1$ & $(xy,yz,zx)$ \\
\bottomrule
\end{tabular}
\end{table}

这张表并不是凭记忆写出的。$T_1$ 行由三维矢量表示的转角迹公式直接得到；二次函数空间的五维无迹部分是 $l=2$ 球谐表示限制到 $O$ 后的结果，它由基函数直接分成 $E\oplus T_2$；剩余两个一维表示再由正交性补齐。加入反演后，$O_h=O\times C_i$，每一行复制为 $g/u$ 两个宇称版本。这就得到 $O_h$ 的十个不可约表示，而不必从头重算十行。

最后还要处理一个在循环群与晶体学表格之间经常造成困惑的问题：复数域上的两个共轭一维表示，在实数域上常合并成一个二维表示。

循环群 $C_n$ 的复不可约表示全是一维，但当 $m$ 与 $n-m$ 不同时，两个表示
\[
\ee^{2\pi\ii m/n},
\qquad
\ee^{-2\pi\ii m/n}
\]
互为复共轭。若物理问题要求使用实基函数，例如 $\cos m\phi$ 与 $\sin m\phi$，这对复一维表示会合并成一个二维实表示。很多化学特征标表中的 $E$ 标签正是在实表示约定下出现的。

这个区别到\chapref{chap:group-theory-quantum}时间反演部分尤其重要：反酉时间反演会把复表示 $D$ 送到 $D^*$。因此两个在纯空间点群下彼此独立的复共轭表示，可能被时间反演配对成简并态。提前区分“复不可约维数”和“实基下合并后的维数”，可以避免把特征标表中的符号机械地当作能级简并度。

\begin{exampleplain}[$O_h$ 中 $d$ 轨道为什么分成 $E_g\oplus T_{2g}$]
五个实 $d$ 轨道可以取为
\[
\{2z^2-x^2-y^2,\ x^2-y^2,\ xy,\ yz,\ zx\}.
\]
它们都是二次齐次函数并且在反演下为偶，因此只能落在 $g$ 型表示中。立方群 $O_h$ 的五维二次谐波表示不是不可约的，而是分解成
\[
\Gamma_{d}=E_g\oplus T_{2g}.
\]
其中
\[
E_g:\quad \{2z^2-x^2-y^2,\ x^2-y^2\},
\]
\[
T_{2g}:\quad \{xy,yz,zx\}.
\]
这个分解不是为了“给轨道贴标签”，而是直接预言：在理想立方晶场中，五重简并的 $d$ 轨道只能劈裂成二重与三重两组。\chapref{chap:group-theory-quantum}会从哈密顿量对称性与微扰理论重新推导这一结论。
\label{ex:oh-d-orbitals}
\end{exampleplain}

对 $D_4$，五个共轭类要求五个不可约表示。四个一维表示由生成元 $r\mapsto\pm1$、$s\mapsto\pm1$ 给出，剩余维数平方和只能由一个二维表示补足。二维自然表示在 $C_4,C_2,C_2',C_2''$ 上的特征标分别为 $0,-2,0,0$，于是得到下表。

\begin{table}[htbp]

\centering
\caption{$D_4$ 的特征标表}
\label{tab:d4-character}
\small
\begin{tabular}{c|ccccc|l}
\toprule
 & $E$ & $2C_4$ & $C_2$ & $2C_2'$ & $2C_2''$ & 基函数 \\
\midrule
$A_1$ & $1$&$1$&$1$&$1$&$1$ & $x^2+y^2,\,z^2$ \\
$A_2$ & $1$&$1$&$1$&$-1$&$-1$ & $R_z$ \\
$B_1$ & $1$&$-1$&$1$&$1$&$-1$ & $x^2-y^2$ \\
$B_2$ & $1$&$-1$&$1$&$-1$&$1$ & $xy$ \\
$E$   & $2$&$0$&$-2$&$0$&$0$ & $(x,y),\,(xz,yz)$ \\
\bottomrule
\end{tabular}

\end{table}



$D_{4h}=D_4\times C_i$，每个表示复制为 $g/u$ 两版，共 10 个不可约表示。

$T_d$ 有五个类，维数平方和给出 $1,1,2,3,3$。其中 $A_2$ 记录不保手性操作的符号，$T_1$ 承载轴矢量，$T_2$ 承载极矢量；$E$ 的特征标由正交关系补齐。这样即可读出表格，而无需把所有行、列正交关系逐对重算。

\begin{table}[htbp]



\centering
\caption{$T_d$ 的特征标表}
\label{tab:td-character}
\begin{tabular}{c|ccccc|l}
\toprule
 & $E$ & $8C_3$ & $3C_2$ & $6S_4$ & $6\sigma_d$ & 基函数 \\
\midrule
$A_1$ & $1$&$1$&$1$&$1$&$1$ & $x^2+y^2+z^2$ \\
$A_2$ & $1$&$1$&$1$&$-1$&$-1$ & $xyz$ \\
$E$   & $2$&$-1$&$2$&$0$&$0$ & $(2z^2-x^2-y^2,\,x^2-y^2)$ \\
$T_1$ & $3$&$0$&$-1$&$1$&$-1$ & $(R_x,R_y,R_z)$ \\
$T_2$ & $3$&$0$&$-1$&$-1$&$1$ & $(x,y,z),\,(xy,yz,zx)$ \\
\bottomrule
\end{tabular}

\end{table}



有了 $T_d$ 的特征标表，转向立方晶系中最重要的 $O_h$。$O_h=O\times C_i$，$O$ 的特征标表已给出（\Cref{tab:o-character}）。每个 $O$ 的不可约表示 $\Gamma$ 产生两个 $O_h$ 表示 $\Gamma_g$（偶宇称）和 $\Gamma_u$（奇宇称），反演 $I$ 在 $g$ 型上取 $+1$，$u$ 型上取 $-1$。

\begin{table}[htbp]



\centering
\caption{$O_h$ 的特征标表（$O\times C_i$，10 个不可约表示）}
\label{tab:oh-character}
\small
\begin{tabular}{c|cccccccccc}
\toprule
 & $E$ & $8C_3$ & $3C_2$ & $6C_4$ & $6C_2'$ & $I$ & $8IC_3$ & $3IC_2$ & $6IC_4$ & $6IC_2'$ \\
\midrule
$A_{1g}$ & $1$&$1$&$1$&$1$&$1$&$1$&$1$&$1$&$1$&$1$ \\
$A_{2g}$ & $1$&$1$&$1$&$-1$&$-1$&$1$&$1$&$1$&$-1$&$-1$ \\
$E_g$   & $2$&$-1$&$2$&$0$&$0$&$2$&$-1$&$2$&$0$&$0$ \\
$T_{1g}$ & $3$&$0$&$-1$&$1$&$-1$&$3$&$0$&$-1$&$1$&$-1$ \\
$T_{2g}$ & $3$&$0$&$-1$&$-1$&$1$&$3$&$0$&$-1$&$-1$&$1$ \\
$A_{1u}$ & $1$&$1$&$1$&$1$&$1$&$-1$&$-1$&$-1$&$-1$&$-1$ \\
$A_{2u}$ & $1$&$1$&$1$&$-1$&$-1$&$-1$&$-1$&$-1$&$1$&$1$ \\
$E_u$   & $2$&$-1$&$2$&$0$&$0$&$-2$&$1$&$-2$&$0$&$0$ \\
$T_{1u}$ & $3$&$0$&$-1$&$1$&$-1$&$-3$&$0$&$1$&$-1$&$1$ \\
$T_{2u}$ & $3$&$0$&$-1$&$-1$&$1$&$-3$&$0$&$1$&$1$&$-1$ \\
\bottomrule
\end{tabular}

\end{table}



$O_h=O\times C_i$，所以 $O$ 的每个不可约表示只需再标记反演宇称。

对 $O$ 的不可约表示 $\Gamma$，偶、奇两种扩展分别满足
\[
\Gamma_g(R)=\Gamma(R),\quad\Gamma_g(IR)=\Gamma(R);
\qquad
\Gamma_u(R)=\Gamma(R),\quad\Gamma_u(IR)=-\Gamma(R).
\]
也就是 $\Gamma_g=\Gamma\otimes g$ 与 $\Gamma_u=\Gamma\otimes u$。

两种扩展彼此正交，并继承 $\Gamma$ 的不可约性；五个 $O$ 表示于是加倍成十个 $O_h$ 表示，维数平方和为 $2(1+1+4+9+9)=48$。

$g$ 型二次谐波给出 $d$ 轨道分裂 $E_g\oplus T_{2g}$，极矢量 $(x,y,z)$ 按 $T_{1u}$ 变换。电偶极跃迁必须满足 $\Gamma_f\otimes T_{1u}\otimes\Gamma_i\supset A_{1g}$，所以同宇称态之间的电偶极跃迁禁戒。
$O_h$ 中常用的张量积分解为
\begin{align*}
E_g\otimes E_g &= A_{1g}\oplus A_{2g}\oplus E_g,\\
T_{1g}\otimes T_{1g} &= A_{1g}\oplus E_g\oplus T_{1g}\oplus T_{2g},\\
E_g\otimes T_{2g} &= T_{1g}\oplus T_{2g},
\end{align*}
它们把允许耦合化成直积中是否含有 $A_{1g}$ 的检验。下面再看两个含反演的轴向群，宇称加倍的机制完全相同。


$D_{3d}=D_3\times C_i$，因此 $D_3$ 的 $A_1,A_2,E$ 各自产生一对 $g/u$ 表示。六个共轭类为
\[
\{E\},\;\{2C_3\},\;\{3C_2\},\;\{i\},\;\{2S_6\},\;\{3\sigma_d\},
\]
其中 $S_6=iC_3$、$\sigma_d=iC_2$。特征标由
\[
\chi_{\Gamma_g}(R)=\chi_{\Gamma}(R),\quad\chi_{\Gamma_g}(iR)=\chi_{\Gamma}(R).
\]
与
\[
\chi_{\Gamma_u}(R)=\chi_{\Gamma}(R),\quad\chi_{\Gamma_u}(iR)=-\chi_{\Gamma}(R).
\]
确定，维数平方和为 $12$。当八面体对称性降到 $D_{3d}$ 时，$T_{2g}$ 限制为
\begin{equation*}
\chi_{T_{2g}}^{O_h}\downarrow_{D_{3d}}=\chi_{E_g}+\chi_{A_{1g}},\qquad \chi_{E_g}^{O_h}\downarrow_{D_{3d}}=\chi_{E_g},
\end{equation*}
所以原来的三重态分成二重态和单态，而 $O_h$ 的 $E_g$ 保持二维。
反演中心还立即推出所有奇阶极张量为零，例如 $D_{3d}$ 晶体没有压电和二阶非线性光学响应；这一结论不需要逐项检查张量分量。

\begin{table}[htbp]



\centering
\caption{$D_{3d}$ 的特征标表}
\label{tab:d3d-character}
\begin{tabular}{c|cccccc}
\toprule
 & $E$ & $2C_3$ & $3C_2$ & $i$ & $2S_6$ & $3\sigma_d$ \\
\midrule
$A_{1g}$ & $1$&$1$&$1$&$1$&$1$&$1$ \\
$A_{2g}$ & $1$&$1$&$-1$&$1$&$1$&$-1$ \\
$E_g$   & $2$&$-1$&$0$&$2$&$-1$&$0$ \\
$A_{1u}$ & $1$&$1$&$1$&$-1$&$-1$&$-1$ \\
$A_{2u}$ & $1$&$1$&$-1$&$-1$&$-1$&$1$ \\
$E_u$   & $2$&$-1$&$0$&$-2$&$1$&$0$ \\
\bottomrule
\end{tabular}

\end{table}



$D_{6h}=D_6\times C_i$，$D_6$ 有 12 元素 6 类，不可约表示维数 $1,1,1,1,2,2$，故 $D_{6h}$ 有 12 个不可约表示。

\begin{table}[htbp]



\centering
\caption{$D_{6h}$ 的特征标表}
\label{tab:d6h-character}
\scriptsize
\setlength{\tabcolsep}{3pt}
\begin{tabular}{c|cccccccccccc}
\toprule
 & $E$ & $2C_6$ & $2C_3$ & $C_2$ & $3C_2'$ & $3C_2''$ & $i$ & $2S_3$ & $2S_6$ & $\sigma_h$ & $3\sigma_d$ & $3\sigma_v$ \\
\midrule
$A_{1g}$ & $1$&$1$&$1$&$1$&$1$&$1$&$1$&$1$&$1$&$1$&$1$&$1$ \\
$A_{2g}$ & $1$&$1$&$1$&$1$&$-1$&$-1$&$1$&$1$&$1$&$1$&$-1$&$-1$ \\
$B_{1g}$ & $1$&$-1$&$1$&$-1$&$1$&$-1$&$1$&$-1$&$1$&$-1$&$1$&$-1$ \\
$B_{2g}$ & $1$&$-1$&$1$&$-1$&$-1$&$1$&$1$&$-1$&$1$&$-1$&$-1$&$1$ \\
$E_{1g}$ & $2$&$1$&$-1$&$-2$&$0$&$0$&$2$&$1$&$-1$&$-2$&$0$&$0$ \\
$E_{2g}$ & $2$&$-1$&$-1$&$2$&$0$&$0$&$2$&$-1$&$-1$&$2$&$0$&$0$ \\
$A_{1u}$ & $1$&$1$&$1$&$1$&$1$&$1$&$-1$&$-1$&$-1$&$-1$&$-1$&$-1$ \\
$A_{2u}$ & $1$&$1$&$1$&$1$&$-1$&$-1$&$-1$&$-1$&$-1$&$-1$&$1$&$1$ \\
$B_{1u}$ & $1$&$-1$&$1$&$-1$&$1$&$-1$&$-1$&$1$&$-1$&$1$&$-1$&$1$ \\
$B_{2u}$ & $1$&$-1$&$1$&$-1$&$-1$&$1$&$-1$&$1$&$-1$&$1$&$1$&$-1$ \\
$E_{1u}$ & $2$&$1$&$-1$&$-2$&$0$&$0$&$-2$&$-1$&$1$&$2$&$0$&$0$ \\
$E_{2u}$ & $2$&$-1$&$-1$&$2$&$0$&$0$&$-2$&$1$&$1$&$-2$&$0$&$0$ \\
\bottomrule
\end{tabular}

\end{table}

$D_6$ 的四个一维表示由 $r\mapsto\pm1$、$s\mapsto\pm1$ 给出；二维表示 $E_1,E_2$ 在主轴转动 $C_6^k$ 上的特征标分别为 $2\cos(k\pi/3)$ 与 $2\cos(2k\pi/3)$，在横向二重轴上均为零。与 $C_i$ 作直积后得到表中的十二个 $D_{6h}$ 表示，维数平方和为 $24$。这张表主要用于六方体系的轨道、声子与跃迁分类；具体能带模型留到后续应用章节。

从特征标表的应用转入衍射物理：空间群对称性不仅约束能级和振动模式，还直接决定衍射谱中的系统消光——这是群论在结构分析中最直接的实验后果。

\begin{theorem}[系统消光定理]



设空间群 $\mathcal G$ 含操作 $\{R|\vb\tau\}$。结构因子
\[
F(\vb G)=\sum_j f_j\,\ee^{-\ii\vb G\cdot\vb r_j}
\]
满足：若对某倒格矢 $\vb G$，$\vb G\cdot\vb\tau\notin2\pi\ZZ$ 且 $R$ 不改变原子类型，则 $F(\vb G)=0$。这给出系统消光规则。

\end{theorem}



结构因子的对称约束如下。若 $\{R|\vb\tau\}\in\mathcal G$ 把原子 $\vb r_j$ 映到同类型原子 $\vb r_{j'}$，则 $f_j=f_{j'}$ 且
\[
\vb r_{j'}=R\vb r_j+\vb\tau.
\]
结构因子
\[
F(\vb G)=\sum_j f_j\,\ee^{-\ii\vb G\cdot\vb r_j}
=\sum_j f_j\,\ee^{-\ii\vb G\cdot(R\vb r_j+\vb\tau)}
\]
（第二个等号是对求和指标做置换 $j\to j'$，因求和遍历所有原子）。提取公共因子：
\[
F(\vb G)=\ee^{-\ii\vb G\cdot\vb\tau}\sum_j f_j\,\ee^{-\ii(R^{\mathsf T}\vb G)\cdot\vb r_j}.
\]

现在考虑 $R^{\mathsf T}\vb G=\vb G$ 的情形。若 $\vb G$ 在 $R$ 下不变（即 $R$ 属于 $\vb G$ 的小群），则
\[
F(\vb G)=\ee^{-\ii\vb G\cdot\vb\tau}F(\vb G).
\]
若 $\vb G\cdot\vb\tau\notin2\pi\ZZ$，则 $\ee^{-\ii\vb G\cdot\vb\tau}\neq1$，故 $F(\vb G)=0$。

把这个判据代入常见的非本原平移即可得到中心化规则。体心平移
$\vb\tau=(\vb a_1+\vb a_2+\vb a_3)/2$ 产生相位 $(-1)^{h+k+l}$，故
$h+k+l$ 为奇数时消光；面心的三个半格矢平移要求 $h+k$、$k+l$、$h+l$
同时为偶数，等价于 $h,k,l$ 全奇或全偶；底心平移
$(\vb a_1+\vb a_2)/2$ 则要求 $h+k$ 为偶数。非对称型操作给出方向相关的条件：
沿 $c$ 轴的 $2_1$ 螺旋把 $(h,k,l)$ 送到 $(-h,-k,l)$，因而只在 $00l$ 方向
保持 $\vb G$ 不变，并使奇数 $l$ 的反射消失；垂直于 $b$ 的 $a$ 滑移面同理使
$h0l$ 中奇数 $h$ 的反射消失。

这些缺峰条件把空间群信息直接编码进衍射图。中心化消光先确定 Bravais 晶格，
螺旋轴和滑移面造成的方向相关消光再区分具有相同点群的空间群。它们来自群作用
强制的结构因子零点，因此与某些原子散射因子偶然抵消造成的零点不同；后者不在
整族反射中稳定出现。反常散射能够破坏非中心对称晶体的 Friedel 等强关系
$I(\vb G)=I(-\vb G)$，但不会解除由空间群强制的系统消光。

\begin{table}[htbp]
\centering
\caption{Bravais 中心化的系统消光规则}
\label{tab:extinction-rules}
\begin{tabular}{c|l}
\toprule
中心化 & 消光条件 \\
\midrule
$P$（原始） & 无系统消光 \\
$I$（体心） & $h+k+l=2n+1$ \\
$F$（面心） & $h,k,l$ 混奇偶 \\
$C$（底心） & $h+k=2n+1$ \\
$R$（菱心） & $-h+k+l=3n\pm1$ \\
\bottomrule
\end{tabular}
\end{table}

系统消光从倒空间给出衍射选择定则。宏观响应的对应结论是 Neumann 原理：
晶体物性张量必须在点群的每个操作下不变。若 $T$ 描述线性响应
$\vb Y=T\vb X$，则
\begin{equation}
TR=RT,
\qquad R\in P.
\label{thm:neumann}
\end{equation}
对 $n$ 阶张量，这个条件写成
\begin{equation}
T_{i_1\cdots i_n}
=\sum_{j_1,\dots,j_n}R_{i_1j_1}\cdots R_{i_nj_n}T_{j_1\cdots j_n}.
\label{eq:neumann-tensor-invariance}
\end{equation}
这些线性约束的解空间就是允许的物性张量空间，其维数给出独立材料常数的数目。

\par\medskip
\noindent\textbf{Neumann 原理的群平均投影}\quad
输入变为 $R\vb X$ 时，输出一方面应为 $T(R\vb X)$，另一方面应为
$R(T\vb X)$，故对任意 $\vb X$ 有
\[
T(R\vb X)=R(T\vb X).
\]
这正是 \eqnrefs{thm:neumann}；在各指标上展开便得到
\eqnrefs{eq:neumann-tensor-invariance}。为直接求这个不变子空间，定义群平均投影
\[
\Pi_T=\frac1{|P|}\sum_{R\in P}R^{\otimes n},
\]
其中 $R^{\otimes n}$ 是 $R$ 在 $n$ 阶张量空间上的诱导作用。则
\[
\Pi_T^2=\frac1{|P|^2}\sum_{R,S}(RS)^{\otimes n}=\frac1{|P|}\sum_{R}R^{\otimes n}=\Pi_T,
\]
故 $\Pi_T$ 是投影。不变张量空间即 $\operatorname{im}\Pi_T$，其维数
\[
\dim\operatorname{im}\Pi_T=\tr\Pi_T=\frac1{|P|}\sum_{R\in P}(\tr R)^n.
\]
对二阶张量（$n=2$），独立分量个数为 $\frac1{|P|}\sum_{R}(\chi_\text{vec}(R))^2$，其中 $\chi_\text{vec}(R)=\tr R$ 是矢量表示的特征标。对称二阶张量（如介电张量 $\epsilon_{ij}$）则需用 $\chi_{\mathrm{Sym}^2}(R)=\frac{\chi_V(R)^2+\chi_V(R^2)}{2}$。

\par\medskip

空间对称约束还需与时间反演约束结合。Onsager 关系在无磁场时给出
$L_{ij}=L_{ji}$，有磁场时给出 $L_{ij}(B)=L_{ji}(-B)$；Neumann 原理则要求同一
输运张量在点群作用下不变。群平均投影因此既可求允许分量，也可检验实验拟合是否
满足晶体对称性。例如一般弹性张量有 21 个独立常数，立方晶体只剩
$C_{11},C_{12},C_{44}$，各向同性介质进一步只剩两个 Lam\'e 常数。

\medskip
\noindent\emph{用 Neumann 原理计数：六方晶系 $D_{6h}$ 的对称二阶张量。}
六方晶系的点群 $D_{6h}=D_6\times C_i$（24 元素）。对称二阶张量在反演下不变（$T_{ij}\to(-1)^2T_{ij}=T_{ij}$），故只需对 $D_6$（12 元素）的纯转动部分求和，结果与 $D_{6h}$ 相同。

$D_6$ 的共轭类、矢量表示特征标 $\chi_V(R)=1+2\cos\theta$（$\theta$ 为绕主轴转角）及 $\chi_{\mathrm{Sym}^2}(R)$：

\[
\begin{array}{c|cccccc}
& E & 2C_6 & 2C_3 & C_2 & 3C_2' & 3C_2'' \\
\hline
\chi_V & 3 & 2 & 0 & -1 & -1 & -1 \\
\chi_V(R^2) & 3 & 0 & 0 & 3 & 3 & 3 \\
\chi_{\mathrm{Sym}^2} & 6 & 2 & 0 & 2 & 2 & 2
\end{array}
\]

其中 $C_6^2=C_3$（$\chi_V=0$），$C_3^2=C_3^{-1}$（$\chi_V=0$），$C_2^2=E$（$\chi_V=3$），$C_2'^2=E$（$\chi_V=3$）。

独立分量数：
\[
N=\frac{1}{12}\bigl[1\cdot6+2\cdot2+2\cdot0+1\cdot2+3\cdot2+3\cdot2\bigr]
=\frac{6+4+0+2+6+6}{12}
=\frac{24}{12}=2.
\]
故六方晶系的对称二阶张量只有 2 个独立分量，对应 $\epsilon=\operatorname{diag}(\epsilon_\perp,\epsilon_\perp,\epsilon_\parallel)$——光轴方向（$c$ 轴）与垂直方向独立，这正是六方晶体光学各向异性的群论来源。

对立方晶系 $O_h$，只需对含 24 个元素的纯转动子群 $O$ 求平均：
\[
\begin{array}{c|ccccc}
& E & 8C_3 & 3C_2 & 6C_4 & 6C_2' \\
\hline
\chi_V & 3 & 0 & -1 & 1 & -1 \\
\chi_V(R^2) & 3 & 0 & 3 & -1 & 3 \\
\chi_{\mathrm{Sym}^2} & 6 & 0 & 2 & 0 & 2
\end{array}
\]
\[
N=\frac{1}{24}[1\cdot6+8\cdot0+3\cdot2+6\cdot0+6\cdot2]=\frac{6+0+6+0+12}{24}=\frac{24}{24}=1.
\]
立方晶系对称二阶张量只有 1 个独立分量 $\epsilon_0 I_3$——光学各向同性。这与已知物理一致：立方晶体在光学上表现为各向同性介质。
\par\medskip

至此，\chapref{chap:point-space-groups}的几何分类与\chapref{chap:representation-theory}的表示论已经真正接上。后续物理不再需要把特征标表视为独立知识点，而是把它当作对称子空间的坐标表。

到这里，本章的几何链条已经闭合：有限旋转群的分类告诉我们点群有哪些；晶体制约定理把它缩成 32 个晶体点群；空间群把点操作与平移统一起来；特征标表则把这些几何对称性转成线性空间中的不可约表示。\chapref{chap:group-theory-quantum}要做的事情，就是把具体量子态、跃迁算符、振动模和 Bloch 波函数放进这些不可约表示中。届时，“对称性决定简并”“某个矩阵元必为零”“某个声子是 Raman 主动”都不再是经验规则，而会从表示的直和与张量积中自动推出。


\subsection{平移不变性与 Bloch 定理}
\label{sec:translation-bloch}


前面的点群都是有限群，而理想晶体还具有无限多个晶格平移。乍看之下，这似乎是完全不同的问题；实际上平移群是 Abel 群，因此它的全部不可约复表示都是一维的。这件事一旦与量子哈密顿量的平移对称性结合，就直接给出 Bloch 定理。

因此 Bloch 定理不必从“猜波函数形状”开始。先取晶格平移 $T_{\vb R}$，其中 $\vb R$ 是任意 Bravais 格矢；由于不同平移可交换，它们可以与 $\hat H$ 同时对角化。一个共同本征态满足
\[
T_{\vb R}|\psi\rangle=\lambda(\vb R)|\psi\rangle.
\]
酉性要求 $|\lambda|=1$，群乘法又要求 $\lambda(\vb R_1+\vb R_2)=\lambda(\vb R_1)\lambda(\vb R_2)$，于是这些一维特征标只能写成 $\ee^{-\ii\vb k\cdot\vb R}$。Bloch 波矢 $\vb k$ 就这样从平移群表示中出现。

前五节主要讨论有限点群。晶体还具有无限离散平移群，因此量子态必须同时尊重平移对称性。Bloch 定理应当从平移群的一维不可约表示推出：它不是某个特定微分方程的巧妙解法，而是 Abel 平移群表示理论的必然结论。


设 Bravais 晶格
\[
\mathcal L
=\{\vb R=n_1\vb a_1+n_2\vb a_2+n_3\vb a_3\}.
\]
对应平移算符定义为
\[
(T_{\vb R}\psi)(\vb r)=\psi(\vb r-\vb R).
\label{eq:translation-operator}
\]
于是
\[
T_{\vb R_1}T_{\vb R_2}=T_{\vb R_1+\vb R_2},
\]
平移群是 Abel 群。

倒格基矢 $\vb b_i$ 由
\[
\vb a_i\cdot\vb b_j=2\pi\delta_{ij}
\]
定义，显式为
\[
\vb b_1=2\pi\frac{\vb a_2\times\vb a_3}{\Omega},
\quad
\vb b_2=2\pi\frac{\vb a_3\times\vb a_1}{\Omega},
\quad
\vb b_3=2\pi\frac{\vb a_1\times\vb a_2}{\Omega},
\]
其中
\[
\Omega=\vb a_1\cdot(\vb a_2\times\vb a_3)
\]
是原胞有向体积。倒格矢写作
\[
\vb G=m_1\vb b_1+m_2\vb b_2+m_3\vb b_3.
\]
由定义立刻有
\[
\ee^{\ii\vb G\cdot\vb R}=1,
\qquad
\forall\vb G\in\mathcal L^*,\ \vb R\in\mathcal L.
\]

无限晶体可以通过 Born--von Karman 周期边界条件先放入有限超胞：
\[
T_{N_i\vb a_i}=I,
\qquad i=1,2,3.
\]
此时平移群成为有限 Abel 群，因此所有不可约表示都是一维。设某个不可约表示在 $T_{\vb a_i}$ 上的值为 $\lambda_i$，则
\[
\lambda_i^{N_i}=1.
\]
于是
\[
\lambda_i=\exp\!\left(-\frac{2\pi\ii m_i}{N_i}\right),
\qquad m_i=0,1,\dots,N_i-1.
\]
把三个整数合并成波矢
\[
\vb k=\sum_i\frac{m_i}{N_i}\vb b_i,
\]
便得到任意平移的不可约表示
\[
D_{\vb k}(T_{\vb R})
=\ee^{-\ii\vb k\cdot\vb R}.
\label{eq:translation-irrep}
\]
由于 $\vb k$ 与 $\vb k+\vb G$ 给出完全相同的表示，波矢只在模倒格矢意义下有区别。因此可以把 $\vb k$ 限制在第一 Brillouin 区中。

\begin{theorem}[Bloch 定理]
若周期哈密顿量满足
\[
[\hat H,T_{\vb R}]=0,
\qquad \forall\vb R\in\mathcal L,
\]
则其本征态可选为平移群不可约表示的共同本征态：
\[
T_{\vb R}\psi_{n\vb k}
=\ee^{-\ii\vb k\cdot\vb R}\psi_{n\vb k}.
\label{eq:bloch-translation-eigenvalue}
\]
按照 \eqnrefs{eq:translation-operator} 的约定，这等价于
\[
\psi_{n\vb k}(\vb r+\vb R)
=\ee^{\ii\vb k\cdot\vb R}
\psi_{n\vb k}(\vb r).
\label{eq:bloch-boundary}
\]
因此波函数可写为
\[
\psi_{n\vb k}(\vb r)
=\ee^{\ii\vb k\cdot\vb r}
 u_{n\vb k}(\vb r),
\qquad
u_{n\vb k}(\vb r+\vb R)
=u_{n\vb k}(\vb r).
\label{eq:bloch-form}
\]
\label{thm:bloch}
\end{theorem}

\par\medskip
\noindent\textbf{Bloch 定理：从平移群一维表示到平面波×周期函数}\quad
先同时对角化。
平移群 $\{T_{\vb R}\}_{\vb R\in\mathcal L}$ 是 Abel 群，且全部 $T_{\vb R}$ 与 $\hat H$ 对易。Abel 群的不可约表示都是一维的，因此可在每个能量本征空间中\emph{同时对角化}所有平移算符。于是可选本征态 $\psi_{n\vb k}$ 满足
\[
T_{\vb R}\psi_{n\vb k}=\ee^{-\ii\vb k\cdot\vb R}\psi_{n\vb k},
\]
其中本征值 $\ee^{-\ii\vb k\cdot\vb R}$ 是平移群的一维不可约表示\eqnrefs{eq:translation-irrep}，标签 $\vb k$ 就是晶体动量。

随后位置表象。
平移算符在位置表象中定义为 $(T_{\vb R}\psi)(\vb r)=\psi(\vb r-\vb R)$，因此本征方程直接给出
\[
\psi_{n\vb k}(\vb r-\vb R)=\ee^{-\ii\vb k\cdot\vb R}\psi_{n\vb k}(\vb r).
\]
把 $\vb r$ 换成 $\vb r+\vb R$，得到更常见的形式
\[
\psi_{n\vb k}(\vb r+\vb R)=\ee^{\ii\vb k\cdot\vb R}\psi_{n\vb k}(\vb r).
\]

接着提取周期因子。
定义
\[
u_{n\vb k}(\vb r)=\ee^{-\ii\vb k\cdot\vb r}\psi_{n\vb k}(\vb r),
\]
则
\begin{align*}
u_{n\vb k}(\vb r+\vb R)
&=\ee^{-\ii\vb k\cdot(\vb r+\vb R)}\psi_{n\vb k}(\vb r+\vb R)\\
&=\ee^{-\ii\vb k\cdot\vb r}\ee^{-\ii\vb k\cdot\vb R}\cdot\ee^{\ii\vb k\cdot\vb R}\psi_{n\vb k}(\vb r)\\
&=\ee^{-\ii\vb k\cdot\vb r}\psi_{n\vb k}(\vb r)=u_{n\vb k}(\vb r),
\end{align*}
即 $u_{n\vb k}$ 具有晶格周期性。因此
\[
\psi_{n\vb k}(\vb r)=\ee^{\ii\vb k\cdot\vb r}\,u_{n\vb k}(\vb r),
\qquad u_{n\vb k}(\vb r+\vb R)=u_{n\vb k}(\vb r),
\]
即"平面波相位 $\times$ 晶格周期函数"。
\par\medskip

Bloch 函数 $\psi_{n\vb k}$ 满足正交归一，能带 $E_n(\vb k)$ 在非简并处连续。由于 $\vb k$ 与 $\vb k+\vb G$ 标记同一个平移不可约表示，只需在第一 Brillouin 区内取代表。时间反演或空间反演对称还会给出 $E_n(\vb k)=E_n(-\vb k)$，进一步减少独立计算区域。

在 $\Gamma$ 点附近，$\vb k\cdot\vb p$ 微扰给出
\begin{align*}
E_n(\vb k)
&=E_n(0)+\frac{\hbar^2k^2}{2m}
+\frac{\hbar^2}{m^2}\sum_{n'\ne n}
\frac{|\langle n0|\vb k\cdot\vb p|n'0\rangle|^2}{E_n(0)-E_{n'}(0)},\\
\frac{1}{m^*_{ij}}
&=\frac{\delta_{ij}}{m}+\frac{2}{m^2}\sum_{n'\ne n}
\frac{\langle n|p_i|n'\rangle\langle n'|p_j|n\rangle}{E_n-E_{n'}}.
\end{align*}
动量算符按 $T_{1u}$ 变换，因此矩阵元非零要求 $\Gamma_n\otimes T_{1u}\otimes\Gamma_{n'}\supset A_{1g}$。这一选择定则把能带的群论约束直接转化为有效质量张量的对称性。

Bloch 基的 Fourier 变换给出 Wannier 函数
\[
w_n(\vb r-\vb R)=\frac{1}{\sqrt N}\sum_{\vb k}
\ee^{-\ii\vb k\cdot\vb R}\psi_{n\vb k}(\vb r).
\]
它们把动量空间的对称信息转化为实空间的定域轨道，是紧束缚模型的自然基底。

Bloch 定理的逻辑核心是：晶格平移群为 Abel 群，其不可约表示是一维相位，而相位标签正是晶体动量 $\vb k$。因此 Bloch 形式不是额外假设，而是平移群不可约表示在位置表象中的必然重写。

由于 $\vb k$ 与 $\vb k+\vb G$ 标记同一个平移不可约表示，我们需要从每个倒格等价类中选一个代表。最对称的选法就是以倒格原点为中心构造 Wigner--Seitz 原胞：对每个倒格点与原点的连线作垂直平分面，取离原点最近的区域。这就是第一 Brillouin 区。

因此 Brillouin 区并不是因为“电子波矢最大到某个值”而人为截断，而是一个商空间基本域：
\[
\mathbb R^3/\mathcal L^*.
\]
边界上的点之所以需要特别小心，是因为它们与相邻 Brillouin 区中的点相差一个倒格矢，常常具有更大的小群或非对称空间群相位。

布里渊区与 Wigner--Seitz 原胞如图 \figref{fig:bz-wigner-seitz} 所示。
\begin{figure}[htbp]
\centering
\begin{tikzfig}[scale=.9]
  \foreach \i in {-2,...,2}{
    \foreach \j in {-2,...,2}{\fill (1.5*\i,1.5*\j) circle (1.1pt);}}
  \draw[line width=.85pt,fill=gray!8] (-.75,-.75) rectangle (.75,.75);
  \node[lecturelabel] at (0,0) {1st BZ};
  \draw[dashed] (-3.2,.75)--(3.2,.75) (-3.2,-.75)--(3.2,-.75)
                (.75,-3.2)--(.75,3.2) (-.75,-3.2)--(-.75,3.2);
\end{tikzfig}
\caption{二维方格倒格子的 Wigner--Seitz 原胞，也就是第一 Brillouin 区。}
\label{fig:bz-wigner-seitz}
\end{figure}

有限周期边界把 $\vb k$ 离散化，而热力学极限把这些点变成 Brillouin 区中的积分测度；两者之间的换算也应在此处一次说明清楚。

Born--von Karman 条件把允许 $\vb k$ 离散化为
\[
\vb k=\sum_i\frac{m_i}{N_i}\vb b_i.
\]
若晶体含 $N=N_1N_2N_3$ 个原胞，第一 Brillouin 区中恰有 $N$ 个不同的 $\vb k$ 点。热力学极限中，$N_i\to\infty$，离散求和变为积分：
\[
\frac1N\sum_{\vb k\in\mathrm{BZ}}f(\vb k)
\longrightarrow
\frac{\Omega}{(2\pi)^3}
\int_{\mathrm{BZ}}\dd^3k\,f(\vb k),
\]
其中 $\Omega$ 为实空间原胞体积。这一测度因子来自倒格原胞体积
\[
\Omega^*=\frac{(2\pi)^3}{\Omega}.
\]
虽然本章重点是群论，但把这一步写清有助于理解为什么能带论里“每个原胞一个 $\vb k$ 态”的计数与平移群表示个数完全一致。

\subsection{Brillouin 区、空间群与小群}
\label{sec:brillouin-lattice-symmetry}

平移对称性告诉我们本征态可以用 $\vb k$ 标记，但第一 Brillouin 区仍包含连续无穷多个点。晶体的点操作还会把不同 $\vb k$ 联系起来，从而进一步减少真正独立的区域。下面从空间群操作 $\{R\mid\vb t\}$ 出发推导 $E_n(R\vb k)=E_n(\vb k)$，并同时引出不可约 Brillouin 区、波矢星与小群。


设空间群元素
\[
g=\{R\mid\boldsymbol\tau\},
\]
其 Hilbert 空间表示记作 $U(g)$。空间群的群关系给出
\[
g\,T_{\vb L}\,g^{-1}=T_{R\vb L}.
\label{eq:space-translation-conjugation}
\]
若 $\ket{\psi_{n\vb k}}$ 是 Bloch 态，满足
\[
T_{\vb L}\ket{\psi_{n\vb k}}
=\ee^{-\ii\vb k\cdot\vb L}
\ket{\psi_{n\vb k}},
\]
则
\begin{align*}
T_{\vb L}U(g)\ket{\psi_{n\vb k}}
&=U(g)T_{R^{-1}\vb L}\ket{\psi_{n\vb k}}\\
&=\ee^{-\ii\vb k\cdot R^{-1}\vb L}
U(g)\ket{\psi_{n\vb k}}\\
&=\ee^{-\ii(R\vb k)\cdot\vb L}
U(g)\ket{\psi_{n\vb k}}.
\end{align*}
因此 $U(g)\ket{\psi_{n\vb k}}$ 属于 $R\vb k$ 的 Bloch 表示。

若 $g$ 是哈密顿量对称操作，便进一步有
\[
E_n(R\vb k)=E_n(\vb k),
\label{eq:band-star-degeneracy}
\]
其中 $R\vb k$ 可随时加上倒格矢折回第一 Brillouin 区。

\begin{definitionplain}[波矢的星与小群]
给定 $\vb k$，其在空间群点操作下得到的所有不等价波矢
\[
\star\vb k
=\{R\vb k\bmod\mathcal L^*\}
\]
称为 $\vb k$ 的\textbf{星}（star of $\vb k$）。

所有满足
\[
R\vb k=\vb k+\vb G,
\qquad \vb G\in\mathcal L^*,
\]
的空间群元素 $\{R\mid\boldsymbol\tau\}$ 构成 $\vb k$ 的\textbf{小群} $\mathcal G_{\vb k}$。把其中纯平移商掉后得到有限的\textbf{小余群}（little co-群）。
\label{def:star-little-group}
\end{definitionplain}

一般 $\vb k$ 点的小群很小，因此能带通常不简并；高对称点、线和面的小群更大，可以出现二维、三维乃至更高维不可约表示，因此常有对称性保护的简并。这正是能带图中 $\Gamma,X,L,K$ 等特殊点需要单独命名的原因。

设 $\vb k_0$ 是高对称点，小群 $\mathcal G_{\vb k_0}$ 较大；沿某条离开它的高对称线，小群通常降为子群 $\mathcal H\subset\mathcal G_{\vb k_0}$。如果端点某能级属于不可约表示 $\Gamma$，那么沿线允许连接到哪些能带，由限制表示
\[
\Gamma\downarrow_{\mathcal H}
\]
决定。这就是能带“兼容关系”的群论来源，与 \secref{sec:perturbative-splitting} 的能级劈裂完全同构：那里是外界微扰把空间对称群降为子群，这里是选定一般波矢本身把保持 $\vb k$ 的对称群降为子群。

如果一个二维不可约表示在沿线小群下分成两个不同的一维表示，那么离开高对称点后两支可以分开；如果沿线仍是二维不可约表示，则二重简并必须沿整条线保留。反之，两条属于不同沿线不可约表示的能带即使能量相遇，也不能杂化打开反交叉；它们可以形成受对称性保护的交叉点。这正是现代能带拓扑讨论中“不同对称本征值可以交叉”的基础版本。

小群还控制一个常用的计数问题：一个 $\vb k$ 点在整个点群作用下会产生多少个对称等价的波矢。

轨道--稳定子定理在倒空间再次出现。若空间群点群为 $P$，而 $P_{\vb k}$ 是保持 $\vb k$ 模倒格矢不变的小余群，则星中不同波矢的个数为
\[
|\star\vb k|=\frac{|P|}{|P_{\vb k}|}.
\label{eq:k-star-size}
\]
因此一般点的小群最小，星最大；高对称点的小群更大，星更小。半导体中常说的“等价谷数”本质上就是某个能谷波矢星的大小。例如若一个最低点位于六个由立方对称联系的方向上，那么六个谷并不是六套独立参数，而是同一星中的六个等价成员。

如果两个波矢属于同一个星，能量由 \eqnrefs{eq:band-star-degeneracy} 联系，因此数值计算没有必要重复求解。把第一 Brillouin 区按点群轨道分割，只保留一个代表区域，就得到\textbf{不可约 Brillouin 区}。

不可约布里渊区如图 \figref{fig:irreducible-bz} 所示。
\begin{figure}[htbp]
\centering
\begin{tikzfig}[scale=1.05]
  \draw[line width=.8pt] (-2.3,-2.3) rectangle (2.3,2.3);
  \draw[fill=gray!12,line width=.7pt] (0,0)--(2.3,0)--(2.3,2.3)--cycle;
  \draw[axes] (-2.7,0)--(2.8,0) node[right] {$k_x$};
  \draw[axes] (0,-2.7)--(0,2.8) node[above] {$k_y$};
  \fill (0,0) circle (1.5pt) node[below left,lecturelabel] {$\Gamma$};
  \fill (2.3,0) circle (1.5pt) node[below right,lecturelabel] {$X$};
  \fill (2.3,2.3) circle (1.5pt) node[above right,lecturelabel] {$M$};
  \node[lecturelabel,fill=white,inner sep=1pt] at (1.7,0.6) {irreducible wedge};
\end{tikzfig}
\caption{具有 $D_4$ 型点对称性的二维方格倒空间。阴影三角形可作为不可约 Brillouin 区；其余区域由旋转和反射得到。}
\label{fig:irreducible-bz}
\end{figure}

如果空间群含有不可消去的分数平移，上面的点群式小群分析还要多保留一个 Bloch 相位。


如果空间群是非对称的，小群表示不能简单退化成点群表示，因为分数平移会带来 Bloch 相位。以二重螺旋轴
\[
S=\{C_{2z}\mid \tfrac c2\hat{\vb z}\}
\]
为例。由 Seitz 乘法
\[
S^2=\{E\mid c\hat{\vb z}\}=T_{c\hat{\vb z}}.
\]
在波矢 $k_z$ 的 Bloch 态上，表示必须满足
\[
D_{\vb k}(S)^2
=\ee^{-\ii k_z c}.
\label{eq:screw-eigenvalue}
\]
因此螺旋轴本征值不是普通点群里的 $\pm1$，而是
\[
s_\pm(k_z)=\pm\ee^{-\ii k_z c/2}.
\]
在 Brillouin 区边界 $k_z=\pi/c$ 时，它们变成 $\pm\ii$。这就是为什么非对称空间群在能带边界上会表现出普通点群表中看不到的结构；再结合时间反演，某些复共轭螺旋本征值还会被迫成对出现。


\subsection{习题与解答}
\label{sec:point-space-exercises}

下面的习题把点群定义、轨道--稳定子、有限旋转群分类、晶格整数结构、晶体制约和 Schoenflies 识别放进具体计算中。解答紧随题目，便于核对每一步使用的定义。

\begin{enumerate}[label=\arabic*.,leftmargin=2.5em]

\item 证明 $SO(3)$ 是 $O(3)$ 的正规子群，并求商群 $O(3)/SO(3)$。

\emph{解答.} 行列式给出满同态 $\det:O(3)\to\{\pm1\}$，其核为 $SO(3)$。故由第一同构定理
\[
O(3)/SO(3)\simeq\{\pm1\}\simeq C_2,
\]
从而 $SO(3)\mathrel{\trianglelefteq} O(3)$。

\item 证明任意 $R\in SO(3)$ 的特征值必可写为 $1,\ee^{\ii\theta},\ee^{-\ii\theta}$。

\emph{解答.} 由 \Cref{thm:rotation-axis}，$1$ 是一个特征值。正交矩阵的复特征值模长为 $1$，且实矩阵的非实特征值成共轭对出现。又因为行列式为 $1$，其余两个特征值乘积为 $1$，故必为 $\ee^{\pm\ii\theta}$。

\item 用双重计数重新推导 \eqnrefs{eq:finite-rotation-counting}。

\emph{解答.} 统计集合
\[
\Omega=\{(g,p)\mid g\in G\setminus\{e\},\ p\in P,\ gp=p\}.
\]
按 $g$ 计数，每个非恒等转动固定两个极点，故 $|\Omega|=2(N-1)$。按极点轨道计数，第 $i$ 个轨道有 $N/n_i$ 个点，每点被 $n_i-1$ 个非恒等稳定子元素固定，故
\[
|\Omega|=\sum_i\frac{N}{n_i}(n_i-1).
\]
令二者相等即得。

\item 由基本方程证明三种多面体旋转群的阶分别为 $12,24,60$。

\emph{解答.} 对 $(2,3,m)$，基本方程给出
\[
\frac12+\frac13+\frac1m=1+\frac2N.
\]
当 $m=3,4,5$ 时分别解得 $N=12,24,60$。

\item 为什么五重旋转可以出现在孤立分子或准晶结构中，却不能成为普通三维周期晶格的点群操作？

\emph{解答.} 点群本身只要求是 $O(3)$ 的有限子群，因此 $I$ 群中的五重轴完全允许。周期晶格还要求旋转在晶格基底下由整数矩阵表示，故其迹必须为整数。五重转动的迹为
\[
1+2\cos\frac{2\pi}{5},
\]
不是整数，违反 \Cref{thm:crystallographic-restriction}，所以普通周期晶格不能有五重轴。

\item 写出 $C_{3v}$ 中矢量表示 $(x,y,z)$ 的不可约分解。

\emph{解答.} 由 \Cref{tab:c3v-character}，$z$ 按 $A_1$ 变换，而 $(x,y)$ 共同按二维表示 $E$ 变换，因此
\[
\Gamma_{\mathrm{vec}}=A_1\oplus E.
\]

\item 写出 $C_{4v}$ 中二次函数空间 $\operatorname{span}\{x^2+y^2,z^2,x^2-y^2,xy,xz,yz\}$ 的不可约分解。

\emph{解答.} 由 \Cref{tab:c4v-character}，$x^2+y^2$ 与 $z^2$ 各承载 $A_1$，$x^2-y^2$ 承载 $B_1$，$xy$ 承载 $B_2$，$(xz,yz)$ 承载 $E$。故
\[
2A_1\oplus B_1\oplus B_2\oplus E.
\]

\item 证明空间群中的平移子群 $T$ 是正规子群。

\emph{解答.} 对任意 $\{R\mid\vb t\}\in\mathcal G$ 和平移 $\{I_3\mid\vb R_0\}$，用 Seitz 乘法得
\begin{align*}
\{R\mid\vb t\}\{I_3\mid\vb R_0\}\{R\mid\vb t\}^{-1}
&=\{I_3\mid R\vb R_0\}.
\end{align*}
由于点操作 $R$ 保持晶格，$R\vb R_0$ 仍是晶格矢量，所以共轭结果仍属于 $T$，故 $T\mathrel{\trianglelefteq}\mathcal G$。

\item 设 $g=\{C_2\mid\vb a/2\}$，其中 $C_2\vb a=\vb a$。求 $g^2$。

\emph{解答.} 由 Seitz 乘法
\[
g^2=\{C_2^2\mid C_2(\vb a/2)+\vb a/2\}
=\{I_3\mid\vb a\}.
\]
这说明二重螺旋轴的平方不是恒等操作，而是一个整格矢平移。

\item 解释 $O_h$ 中下标 $g/u$ 的来源。

\emph{解答.} 因为 $O_h=O\times C_i$，而 $C_i$ 有两个一维不可约表示，分别把反演 $\mathcal I$ 映到 $+1$ 和 $-1$。将 $O$ 的任一不可约表示与这两个一维表示做张量积，就得到偶宇称 $g$ 与奇宇称 $u$ 两个版本。

\item 利用维数平方和验证 $I$ 群不可约表示维数 $1,3,3,4,5$ 与群阶 $60$ 相容。

\emph{解答.}
\[
1^2+3^2+3^2+4^2+5^2=1+9+9+16+25=60.
\]

\item 对正四面体完整点群 $T_d$，说明为什么它不是 $T\times C_i$。

\emph{解答.} $T_d$ 不包含空间反演 $\mathcal I$，其不保手性操作主要由镜面和 $S_4$ 型操作组成，因此不存在中央的独立 $C_i$ 因子。包含反演并与 $T$ 构成直积的是 $T_h=T\times C_i$，不是 $T_d$。

\end{enumerate}
